% ===================================================================
%  LENTELĖ Nr. 16 — Didžioji parduotuvė. — Bazaras. Žaislai.
%  Traduction 2026 d'apres texto/io, controlee sur texto/fr, texto/ru
%  et texto/cs. Consigne : voir texto/lt/10-lentele-01.tex.
%
%  DERNIER TABLEAU DE LA COLONNE, et la menagerie en carton de
%  l'alinea 3 lui donne son sujet. LE GALICIEN AVAIT TROUVE LA
%  MENAGERIE PARMI LES ENDROITS OU LA DISTANCE A LA LANGUE VOISINE
%  EST MAXIMALE ; le lituanien y montre autre chose, et de plus
%  precis : LA LIGNE ENTRE CE QU'IL NOMME ET CE QU'IL EMPRUNTE Y
%  PASSE ENTRE DEUX MANIERES DE CONNAITRE UNE BETE.
%
%  Portent un nom lituanien, et tous sont des composes ou des
%  derives transparents :
%
%    « šikšnosparnis » la chauve-souris, l'aile-de-cuir ;
%    « ryklys »       le requin, de « ryti » devorer ;
%    « banginis »     la baleine ;
%    « vėžlys »       la tortue ;
%    « ruonis »       le phoque ;
%    « šiaurės elnias » le renne, le cerf du nord ;
%    « beždžionė »    le singe ;
%    et bien sur « vilkas », « lapė », « šeškas », « žiurkė »,
%    « pelė », qui vivent dans le pays.
%
%  Restent etrangers : « rinoceras », « begemotas », « krokodilas »,
%  « leopardas », « pantera », « hiena », « strutis », « papūga »,
%  « dromedaras », « boa ».
%
%  ET LE CHAMEAU TRAVERSE LA LIGNE. « kupranugaris » est un compose
%  lituanien — « kupra » la bosse, « nugara » le dos, le
%  bosse-au-dos — pour une bete qu'aucun Lituanien n'a jamais vue
%  chez lui. Pourquoi celle-la et non le rhinoceros ? PARCE QU'ELLE
%  EST DANS L'EVANGILE. Le chameau qui passe par le chas d'une
%  aiguille a du etre nomme avant qu'aucun zoo n'arrive, et c'est le
%  livre qui l'a nomme.
%
%  LE CRITERE DU LIEU D'APPRENTISSAGE GAGNE DONC ICI SA DERNIERE
%  PRECISION, ET C'EST LA PLUS ETROITE : le lieu peut etre un LIVRE,
%  et un livre en particulier. Le tableau 14 avait oppose l'atelier
%  et le diplome ; le tableau 15, la date d'arrivee d'une plante ; le
%  16 ajoute qu'un mot peut entrer dans une langue par un seul texte,
%  des siecles avant la chose.
%
%  L'ECART DECLARE DU BLOC t16-15-1, LE TROISIEME ET DERNIER DES
%  TROIS. Le fac-simile ido y ecrit « mi-sferi d) », sans la
%  parenthese ouvrante. Comme au tableau 8, la coquille est du
%  compositeur de l'ido et non du texte lituanien : la parenthese est
%  retablie, comme dans les treize colonnes precedentes, et
%  renvoji.py passe l'ecart sans rien dire.
% ===================================================================

%%K t16-apar-1 apar
\VUsaut{4.0mm}
\VUtitre{15.0pt}{60}{LENTELĖ Nr. 16}
\VUsaut{2.0mm}
\VUpk{13.2pt}{40}{Didžioji parduotuvė. --- Bazaras.}
\VUpk{13.2pt}{40}{Žaislai.}
\VUsaut{3.4mm}

%%K t16-01-1 p
1. --- Artėja Nauji metai. Didžiosios parduotuvės parengė savo \VUgras{žaislų}
\textsuperscript{(1)} ir naujametinių dovanų vitrinas.

%%K t16-01-2 p
Paprašiau senelio, kad nuvestų mane į bazarą pamatyti tos gražios parodos, ir jis
sutiko.

%%K t16-02-1 p
2. --- Pirmiausia sustojome prieš gražius ginklų skydus, kuriuose yra
\VUgras{šautuvas}, \VUgras{kepurė}, \VUgras{kardas}, \VUgras{špagos diržas} ir
\VUgras{antpečių} pora, arba \VUgras{šalmas}, \VUgras{šarvas} \textsuperscript{(3)} ir
\VUgras{pistoletas} \textsuperscript{(4)}. Visa tai mane sudomino daug labiau negu
\VUgras{šviesos raketos} \textsuperscript{(2)}.

%%K t16-tit-1 sub
\VUsaut{3.0mm}
\VUpk{10.2pt}{40}{Gražūs reginiai.}
\VUsaut{2.6mm}

%%K t16-03-1 p
3. --- Tame pačiame skyriuje buvo \VUgras{sargybinis} \textsuperscript{(5)} savo
\VUgras{būdelėje}; jis tarsi ėjo sargybą prieš visą kartoninį žvėryną. Kiek ten buvo
gyvūnų: \VUgras{Papūga} \textsuperscript{(6)} ant savo laktos, \VUgras{raganosis}
\textsuperscript{(7)} su ragu ant nosies, \VUgras{šiaurės elnias} \textsuperscript{(8)},
\VUgras{strutis} \textsuperscript{(9)}, \VUgras{beždžionė} \textsuperscript{(10)},
gliaudanti riešutą, \VUgras{lapė} \textsuperscript{(11)}, nešanti vištą,
\VUgras{krokodilas} \textsuperscript{(12)}, žiojantis milžiniškais
\VUgras{žandikauliais}, \VUgras{kupranugaris} \textsuperscript{(13)}, elegantiškas
\VUgras{kurtas} \textsuperscript{(14)}, \VUgras{pantera} \textsuperscript{(15)}, paskui du
gyvūnai, kurių nepažinau ir kurie, kaip sakoma, yra \VUgras{šeškas}
\textsuperscript{(16)} ir \VUgras{hiena} \textsuperscript{(17)}. Taip pat ten buvo
\VUgras{leopardas} \textsuperscript{(18)} ir \VUgras{vilkas} \textsuperscript{(21)};
visai netoli buvo dailūs \VUgras{žiurkiukai} \textsuperscript{(19)} ir mažytės
\VUgras{pelytės} \textsuperscript{(20)} iš gumos. Galiausiai \VUgras{begemotas}
\textsuperscript{(22)} ir kuprotas \VUgras{dromedaras} \textsuperscript{(23)} užbaigė
rinkinį. Būčiau ten likęs dar daugelį valandų, jei šalia nebūtų buvusi puiki
stovykla, pilna \VUgras{švininių kareivėlių} \textsuperscript{(29)}, kuri patraukė
mano dėmesį.

%%K t16-tit-2 sub
\VUsaut{4.0mm}
\VUpk{10.2pt}{40}{Kareiviai ir karas.}
\VUsaut{2.8mm}

%%K t16-04-1 p
4. --- Mano senelis, kuris yra \VUgras{invalidas} \textsuperscript{(24)} ir kuriam
teko palikti kariuomenę dėl \VUgras{žaizdos,} pelniusios jam \VUgras{karo medalį,}
labai mėgsta pasakoti man \VUgras{kampanijas,} kuriose dalyvavo. Jis buvo laimingas
aiškindamas įvairius mažosios miniatiūrinės kariuomenės judesius. Būrelis tikrų
kareivių, lankančių bazarą --- \VUgras{inžinerijos kareivis} \textsuperscript{(25)},
\VUgras{Alpių šaulys} \textsuperscript{(26)} su \VUgras{berete} ant galvos, pėstininkų
\VUgras{viršila} \textsuperscript{(27)} ir \VUgras{artilerijos seržantas}
\textsuperscript{(28)} su auksiniu \VUgras{galionu} ant rankovės ---, sustojo prie mūsų
paklausyti paaiškinimų.

%%K t16-05-1 p
5. --- Mano senelis, labai išdidus, kad turi tokių puikių klausytojų,
improvizavo ištisą mūšio planą.

%%K t16-05-2 p
Įtvirtintas miestas su savo \VUgras{pylimais} ir \VUgras{grioviais} buvo apgultas
\VUgras{priešo kariuomenės,} kuri, po ilgo \VUgras{bombardavimo,} buvo pasirengusi
\VUgras{šturmuoti.} Bet \VUgras{apsiaustieji,} bado priversti, mėgino
\VUgras{pulti iš apsupties} į \VUgras{apgulėjų} \VUgras{stovyklą.} Atrėmę
\VUgras{puolimą} atkakliomis \VUgras{kautynėmis} aplink \VUgras{palapines,}
\VUgras{nugalėję} apgulėjai paleido savo \VUgras{eskadronus} persekioti
\VUgras{nugalėtų} užpuolikų. Šie \VUgras{traukėsi,} nuklodami \VUgras{kautynių lauką}
\VUgras{žuvusiais} ir \VUgras{sužeistais.} \VUgras{Sanitarai} nunešė pastaruosius į
\VUgras{ligoninę,} kad juos slaugytų \VUgras{chirurgas.}

%%K t16-05-3 p
Nepaisant kelių \VUgras{batalionų,} sudariusių \VUgras{ketvirtainį,} didvyriško
pasipriešinimo, \VUgras{pralaimėjimas} buvo visiškas, kariuomenės likutis
\VUgras{pabėgo,} palikdamas daug \VUgras{belaisvių.} Priešas nuėjo užbaigti savo
\VUgras{pergalės,} užimdamas miestą. Galbūt tai bus kampanijos pabaiga, ir po
\VUgras{paliaubų} bus galima pasirašyti \VUgras{taikos sutartį.}

%%K t16-tit-3 sub
\VUsaut{4.4mm}
\VUpk{10.2pt}{40}{Naujametinės dovanos.}
\VUsaut{2.8mm}

%%K t16-06-1 p
6. --- Jaunieji kareiviai, kurie juokėsi klausydami vaizdingo veterano pasakojimo,
paprašė jo dar kelių paaiškinimų. O aš apėjau prekystalį pažiūrėti gražaus
\VUgras{arbaleto} \textsuperscript{(32)} ir \VUgras{lanko} \textsuperscript{(33)} su jo
\VUgras{strėle} \textsuperscript{(31)}. Ten radau kitą gyvūnų rinkinį: du
\VUgras{giedančius paukščius} \textsuperscript{(34)} --- automatinę
\VUgras{lakštingalą} ir \VUgras{devynbalsę,} giedančią visa gerkle --- ir keistų
gyvūnų eilę: \VUgras{šikšnosparnį} \textsuperscript{(35)}, \VUgras{smauglį}
\textsuperscript{(36)}, \VUgras{vėžlį} \textsuperscript{(37)}, \VUgras{ruonį}
\textsuperscript{(38)}, \VUgras{banginį} \textsuperscript{(39)} ir pripučiamą
\VUgras{ryklį} \textsuperscript{(40)}.

%%K t16-07-1 p
7. --- Neilgai stovėjau į juos žiūrėdamas, nes pamačiau tai, apie ką svajojau: puikų
\VUgras{marionečių teatrą} \textsuperscript{(41)} su nuostabiomis
\VUgras{dekoracijomis} ir žavinga raudona \VUgras{uždanga.} \VUgras{Scenoje} du
aktoriai tarsi vaidino \VUgras{vaidmenį} labai įdomiame \VUgras{veikale,} nes mažieji
kartoniniai \VUgras{žiūrovai,} įtaisyti ložėse arba sėdintys \VUgras{krėsluose} ir
\VUgras{parteryje} už \VUgras{orkestro vadovo,} gyvai plojo.

%%K t16-07-2 p
Gaila, tas teatras kainavo labai brangiai.

%%K t16-08-1 p
8. --- Beje, išsirinkti žaislus buvo sunku, nes vitrinose buvo sukrauta visokiausių
žaislų: \VUgras{Arlekinas} \textsuperscript{(42)}, \VUgras{Pulčinelė}
\textsuperscript{(43)}, \VUgras{Pjero} \textsuperscript{(44)}, \VUgras{apuokai}
\textsuperscript{(45)}, plasnojantys sparnais, švilpiantys \VUgras{juodieji strazdai}
\textsuperscript{(46)}, lojantys \VUgras{pudeliai}, \VUgras{fonografas}
\textsuperscript{(47)} ir \VUgras{kantrybės žaislai.} Vienas iš tų žaislų mane labai
palinksmino: jis vaizdavo \VUgras{jūrų mūšį} \textsuperscript{(48)}, kuriame
\VUgras{laivą} puolė \VUgras{piratai} prie krašto, apsodinto \VUgras{kokosų
palmėmis.}

%%K t16-noto-1 noto
\VUnotes{143.2mm}{%
\textit{(*) Žiūrėk lentelę Nr.~6.}%
}

%%K t16-09-1 p
9. --- Galėjau labai lengvai apžiūrėti visus tuos stebuklus, nes parduotuvėje buvo
nedaug žmonių, vos keli dykinėtojai, \VUgras{dragūnas} \textsuperscript{(49)} ir
\VUgras{inkasatorius} \textsuperscript{(50)}, kuris pateikė \VUgras{vekselį}
pagrindinėje \VUgras{kasoje} \textsuperscript{(51)}.

%%K t16-10-1 p
10. --- Paklausiau senelio, kam skirtos knygos, sudėtos ant lentynų už
\VUgras{kasininkės}; jis atsakė man, kad tai \VUgras{apskaitos knygos.}

%%K t16-tit-4 sub
\VUsaut{3.6mm}
\VUpk{10.2pt}{40}{Žiopsojimas prie krautuvės.}
\VUsaut{2.8mm}

%%K t16-11-1 p
11. --- Palikę žaislų skyrių, nuėjome į balnų skyrių dirstelėti į \VUgras{balnus}
\textsuperscript{(53)}, \VUgras{balnakilpes} \textsuperscript{(54)}, \VUgras{kamanas}
\textsuperscript{(55)}, \VUgras{žąslus} \textsuperscript{(56)} ir \VUgras{botagus.} Kol
\VUgras{skyriaus vedėjas} \textsuperscript{(57)} teikė kelias žinias mano seneliui, aš
nuėjau pažiūrėti \VUgras{porceliano} ir \VUgras{krištolo} \textsuperscript{(58)}
vitrinos, kurioje yra gražių \VUgras{salotinių} ir dailių \VUgras{arbatinukų}
\textsuperscript{(59)}.

%%K t16-12-1 p
12. --- Norėdami eiti į antrąjį aukštą, praėjome už \VUgras{korsetų}
\textsuperscript{(60)} vitrinos, šalia kojinių skyriaus, kur buvo išdėtos labai gražios
\VUgras{kelnaitės} \textsuperscript{(63)}. Netoliese \VUgras{pardavėja}
\textsuperscript{(61)} leido \VUgras{pirkėjai} \textsuperscript{(62)} rinktis gražius
šilko \VUgras{audinius} \textsuperscript{(64)}.

%%K t16-12-2 p
Lipdamas laiptais pastebėjau, kad parduotuvė papuošta japoniškais
\VUgras{žibintais} \textsuperscript{(65)}, \VUgras{skėčiais} \textsuperscript{(66)} ir
milžiniškomis \VUgras{palmėmis} \textsuperscript{(70)}.

%%K t16-13-1 p
13. --- Antrajame aukšte nebuvo daug ko matyti; tik baldai (*), \VUgras{seifai}
\textsuperscript{(67)}, \VUgras{komodos} \textsuperscript{(68)}, \VUgras{vaikiški
vežimėliai} \textsuperscript{(69)}. Bet bendras parduotuvės vaizdas buvo labai
\VUgras{linksmas.}

%%K t16-14-1 p
14. --- Prie mūsų kojų mačiau muzikos instrumentus: \VUgras{arfas}
\textsuperscript{(71)}, \VUgras{gitaras} \textsuperscript{(72)}, \VUgras{mandolinas}
\textsuperscript{(73)}, \VUgras{korneta-pistonus} \textsuperscript{(74)},
\VUgras{klarnetus} \textsuperscript{(75)}, smuikus su jų \VUgras{smičiumi}
\textsuperscript{(76)} ir net \VUgras{didįjį būgną} \textsuperscript{(77)}. Kiek
toliau, prie popieriaus skyriaus, \VUgras{studentas} \textsuperscript{(78)} derėjosi su
\VUgras{pardavėju} \textsuperscript{(79)} dėl sąsiuvinių dėklo. Šalia jų įžiūrėjau
gražų \VUgras{elementorių} \textsuperscript{(80)}.

%%K t16-14-2 p
Parduotuvės viduryje buvo pastatyta aukšta \VUgras{Kalėdų eglutė}
\textsuperscript{(81)}, visa apkarstyta žvakelėmis, niekučiais ir visokiais žaislais:
\VUgras{mediniais arkliukais} \textsuperscript{(82)}, \VUgras{lėlėmis}
\textsuperscript{(83)} ir t.~t.

%%K t16-14-3 p
\VUgras{Kvepalų skyrius} \textsuperscript{(84)} su savo \VUgras{dantų milteliais} ir
įvairiais \VUgras{kvepalais} skleidė labai gerą kvapą, kuris sklido iki mūsų.

%%K t16-tit-5 sub
\VUsaut{3.4mm}
\VUpk{10.2pt}{40}{Šį tą perkame.}
\VUsaut{2.8mm}

%%K t16-15-1 p
15. --- Kadangi mano senelis ėmė laikyti laiką per ilgu, nusileidome didžiaisiais
laiptais. Jis kiek stabtelėjo prieš \VUgras{astronomijos lentelę}
\textsuperscript{(85)}, vaizduojančią, kaip jis man pasakė, \VUgras{kometas}
\textsuperscript{(a)}, \VUgras{mėnulio ir saulės užtemimus} \textsuperscript{(b)},
vėjų rožę su \VUgras{keturiomis pasaulio šalimis} \textsuperscript{(c)}
\VUgras{(Šiaure, Pietumis, Rytais ir Vakarais)}, dviejų \VUgras{pusrutulių}
\textsuperscript{(d)} žemėlapį, \VUgras{planetas} \textsuperscript{(e)} ir
\VUgras{saulės sistemą} \textsuperscript{(f)}. Nieko iš viso to nesupratau, ir mane
daug labiau domino galionuota \VUgras{išvežiotojo} \textsuperscript{(86)} livrėja,
kuris praėjo, ir gražios \VUgras{vėduoklės} \textsuperscript{(87)}, kurios buvo šalia
manęs, prie \VUgras{piniginių} \textsuperscript{(88)} ir \VUgras{portfelių}
\textsuperscript{(89)}.

%%K t16-16-1 p
16. --- Taip pat gėrėjausi prekystalyje \VUgras{plunksnomis} \textsuperscript{(90)} ir
\VUgras{diržų sagtimis} \textsuperscript{(91)}, gražiomis \VUgras{dirbtinėmis
gėlėmis} \textsuperscript{(94)}, grakščiai sudėtomis į krepšius. \VUgras{Žibuoklės},
\VUgras{leukojai}, \VUgras{hiacintai} \textsuperscript{(95)}, \VUgras{pelargonijos}
\textsuperscript{(96)}, \VUgras{lelijos} \textsuperscript{(97)} ir \VUgras{nasturtos}
\textsuperscript{(98)} buvo labai gerai nukopijuotos.

%%K t16-tit-6 sub
\VUsaut{4.4mm}
\VUpk{10.2pt}{40}{Senelio dovana.}
\VUsaut{2.8mm}

%%K t16-17-1 p
17. --- Paprašiau senelio nupirkti man vieną iš \VUgras{dantų šepetėlių}
\textsuperscript{(92)}, kurie buvo dėžutėje šalia \VUgras{plaukų smeigtukų}
\textsuperscript{(93)}. Jis sutiko, ir aš pats nuėjau sumokėti už savo pirkinį prie
kasos, prie kurios buvo ruošiama svarbi \VUgras{siunta} \textsuperscript{(100)}
\VUgras{klientui} \textsuperscript{(99)}.

%%K t16-18-1 p
18. --- Staiga lengvas sujudimas kilo tarp žmonių, sustojusių prie meninių
\VUgras{bronzos dirbinių} \textsuperscript{(101)}. \VUgras{Vagis}
\textsuperscript{(102)} ką tik buvo sučiuptas nusikaltimo vietoje
\VUgras{inspektoriaus} \textsuperscript{(103)}, kai mėgino ką nors apiplėšti.
Pasirūpinta atvesdinti policininką, kad jį \VUgras{suimtų,} ir kaip tik kai išėjome,
nusikaltėlis buvo vedamas į kalėjimą.

\VUsaut{6.0mm}
\VUfilet{22mm}
